\documentclass[aspectratio=169,11pt]{beamer}
\usepackage{../phanes-beamer}
\renewcommand{\ModuleID}{06}
\title{Operating local inference}
\subtitle{Size, serve, benchmark, and maintain models from 27B to 2.4T}
\author{PHANES Initiative}
\institute{The University of Texas at Austin}
\date{September 12, 2026}
\begin{document}
\begin{frame}[plain]
\titlepage
\end{frame}
\begin{frame}[t,fragile]{Start with the working NVFP4 launcher}
\fontsize{11.5}{14}\selectfont
\begin{itemize}\setlength{\itemsep}{2mm}
\item Primary: vllm/vllm-qwen38-nvfp4.sh.
\item Served alias: qwen3.8-27b-nvfp4; HTTPS port 41883.
\item The supplied baseline already uses MTP with three draft tokens.
\item DFlash2 remains a separate experimental extension.
\end{itemize}
\vfill{\fontsize{6.5}{8}\selectfont\color{UTgray} Sources: user. See references and instructor notes.\par}
\end{frame}
\begin{frame}[t,fragile]{A serving budget has several components}
\fontsize{11.5}{14}\selectfont
\begin{itemize}\setlength{\itemsep}{2mm}
\item Weights, including quantization metadata and unquantized tensors.
\item KV cache and any recurrent or hybrid model state.
\item Activations, workspaces, graph captures, and runtime allocations.
\item Headroom for long prefill, concurrency, and operational variation.
\end{itemize}
\end{frame}
\begin{frame}[t,fragile]{Estimate weights with an explicit lower bound}
\fontsize{11.5}{14}\selectfont
\begin{itemize}\setlength{\itemsep}{2mm}
\item Raw weight bytes $\approx$ total parameters $\times$ bits per parameter / 8.
\item Decimal GB = bytes / $10^9$; $\text{GiB} = \text{bytes} / 2^{30}$.
\item Use total stored parameters when sizing a mixture-of-experts model.
\item Add actual checkpoint overhead and measured runtime needs.
\end{itemize}
\end{frame}
\begin{frame}[t,fragile]{Weight-only memory grows quickly}
\fontsize{10.5}{12.5}\selectfont\setlength{\tabcolsep}{4pt}\renewcommand{\arraystretch}{1.15}
\begin{tabular}{>{\raggedright\arraybackslash}p{\dimexpr 0.4\linewidth-12.8pt\relax}>{\raggedright\arraybackslash}p{\dimexpr 0.2\linewidth-6.4pt\relax}>{\raggedright\arraybackslash}p{\dimexpr 0.2\linewidth-6.4pt\relax}>{\raggedright\arraybackslash}p{\dimexpr 0.2\linewidth-6.4pt\relax}}
\toprule
\textbf{Total parameters} & \textbf{16-bit GB} & \textbf{8-bit GB} & \textbf{4-bit GB} \\ \midrule
27B & 54 & 27 & 13.5 \\[2pt]
70B (generic tier) & 140 & 70 & 35 \\[2pt]
320B (GLM rounded) & 640 & 320 & 160 \\[2pt]
2.4T (Qwen Max) & 4,800 & 2,400 & 1,200 \\[2pt]
\bottomrule\end{tabular}\par\vspace{2mm}
\vfill{\fontsize{6.5}{8}\selectfont\color{UTgray} Sources: glm, qmax. See references and instructor notes.\par}
\end{frame}
\begin{frame}[t,fragile]{Use model size and workload together}
\fontsize{10.5}{12.5}\selectfont\setlength{\tabcolsep}{4pt}\renewcommand{\arraystretch}{1.15}
\begin{tabular}{>{\raggedright\arraybackslash}p{\dimexpr 0.25\linewidth-6.0pt\relax}>{\raggedright\arraybackslash}p{\dimexpr 0.35\linewidth-8.4pt\relax}>{\raggedright\arraybackslash}p{\dimexpr 0.4\linewidth-9.6pt\relax}}
\toprule
\textbf{Model case} & \textbf{Architecture fact} & \textbf{Teaching role} \\ \midrule
Qwen3.8-27B NVFP4 & 27B class; supplied launch script & Primary class endpoint \\[2pt]
GLM-5.3-Flash & $\approx 320$B total / 18B active & Large multimodal MoE comparison \\[2pt]
Qwen3.8-Max & 2.4T total / 95B active & Upper-end infrastructure sizing \\[2pt]
\bottomrule\end{tabular}\par\vspace{2mm}
\vfill{\fontsize{6.5}{8}\selectfont\color{UTgray} Sources: user, glm, qmax. See references and instructor notes.\par}
\end{frame}
\begin{frame}[t,fragile]{Why 18B active does not mean an 18B fit}
\fontsize{11.5}{14}\selectfont
\begin{itemize}\setlength{\itemsep}{2mm}
\item GLM-5.3-Flash activates a subset of its experts per token.
\item Its complete checkpoint still contains roughly 320B parameters.
\item Compare raw FP8 weights with one 96 GB GPU.
\item What changes if experts are kept in host RAM?
\end{itemize}
\vfill{\fontsize{6.5}{8}\selectfont\color{UTgray} Sources: glm, rtx. See references and instructor notes.\par}
\end{frame}
\begin{frame}[t,fragile]{Worked solution: why 18B active does not fit 96 GB}
\fontsize{11.5}{14}\selectfont
\begin{itemize}\setlength{\itemsep}{2mm}
\item Weight bytes: about 320B parameters $\times$ 8 bits / 8 $\approx$ 320 GB in FP8, which exceeds one 96 GB GPU by a factor of about 3.3.
\item "18B active" describes only the expert subset applied per token; the checkpoint still stores the full roughly 320B parameters.
\item Keeping experts in host RAM moves where the bytes reside but does not remove the checkpoint; it adds transfer cost and requires explicit runtime support.
\item The decision is therefore a multi-GPU or multi-node layout, checked against measured usable bytes, not a choice of which 18B subset to keep.
\end{itemize}
\vfill{\fontsize{6.5}{8}\selectfont\color{UTgray} Sources: glm, rtx. See references and instructor notes.\par}
\end{frame}
\begin{frame}[t,fragile]{Consider three hardware classes}
\fontsize{10.5}{12.5}\selectfont\setlength{\tabcolsep}{4pt}\renewcommand{\arraystretch}{1.15}
\begin{tabular}{>{\raggedright\arraybackslash}p{\dimexpr 0.25\linewidth-6.0pt\relax}>{\raggedright\arraybackslash}p{\dimexpr 0.35\linewidth-8.4pt\relax}>{\raggedright\arraybackslash}p{\dimexpr 0.4\linewidth-9.6pt\relax}}
\toprule
\textbf{Class} & \textbf{Capacity illustration} & \textbf{What to investigate} \\ \midrule
One 96 GB GPU & 96 GB nominal VRAM & 27B baseline; profile long context \\[2pt]
Four 96 GB GPUs & 384 GB aggregate & GLM FP8 fit margin and PCIe cost \\[2pt]
Eight 288 GB GPUs & 2,304 GB aggregate & Trillion-scale quantized weights \\[2pt]
\bottomrule\end{tabular}\par\vspace{2mm}
\vfill{\fontsize{6.5}{8}\selectfont\color{UTgray} Sources: rtx, b300. See references and instructor notes.\par}
\end{frame}
\begin{frame}[t,fragile]{Size GLM for four 96 GB GPUs}
\fontsize{11.5}{14}\selectfont
\begin{itemize}\setlength{\itemsep}{2mm}
\item Nominal aggregate memory: 384 GB.
\item Idealized 320B FP8 weights: 320 GB, leaving 64 GB.
\item The recipe reports about 306 GiB of checkpoint weights.
\item What else must fit, and what evidence decides feasibility?
\end{itemize}
\vfill{\fontsize{6.5}{8}\selectfont\color{UTgray} Sources: glmrecipe. See references and instructor notes.\par}
\end{frame}
\begin{frame}[t,fragile]{Worked solution: the GLM sizing gap}
\fontsize{11.5}{14}\selectfont
\begin{itemize}\setlength{\itemsep}{2mm}
\item The recipe checkpoint is about 306 GiB, roughly 328.6 decimal GB, so the weights leave about 55.4 GB of the nominal 384 GB.
\item What else must fit: the KV cache for live requests, activations and workspace, runtime overhead, and the per-device split of a four-way partition.
\item Feasibility evidence: each GPU's actual usable bytes, a measured memory profile at representative context and concurrency, and kernel support for the precision on this hardware.
\item Until that evidence exists, 55.4 GB is a nominal margin, not a guarantee; do not promise the recipe's full 1M context.
\end{itemize}
\vfill{\fontsize{6.5}{8}\selectfont\color{UTgray} Sources: glmrecipe. See references and instructor notes.\par}
\end{frame}
\begin{frame}[t,fragile]{Size Qwen Max beyond a workstation}
\fontsize{11.5}{14}\selectfont
\begin{itemize}\setlength{\itemsep}{2mm}
\item Idealized 4-bit weights: 1.2 TB; 8-bit weights: 2.4 TB.
\item An eight-B300 example provides 2.304 TB nominal GPU memory.
\item 8-bit weights alone exceed that capacity.
\item 4-bit arithmetic leaves room, but is not a deployment guarantee.
\end{itemize}
\vfill{\fontsize{6.5}{8}\selectfont\color{UTgray} Sources: qmax, b300. See references and instructor notes.\par}
\end{frame}
\begin{frame}[t,fragile]{Not every served model is a chat model}
\fontsize{11.5}{14}\selectfont
\begin{itemize}\setlength{\itemsep}{2mm}
\item An embedding model turns text into a fixed-size vector instead of generating tokens; it has no chat loop and no tool calls.
\item Retrieval quality depends on the embedding revision, so an index must record the model version that built it.
\item Embedding workloads are short and batch-oriented, so a smaller model often serves them well at low memory cost.
\item The class LLM endpoint does not automatically provide embeddings; serving one is a separate configuration with its own budget and contract.
\end{itemize}
\end{frame}
\begin{frame}[t,fragile]{Context can dominate the remaining memory}
\fontsize{11.5}{14}\selectfont
\begin{itemize}\setlength{\itemsep}{2mm}
\item A simple attention KV estimate depends on layers, KV heads, and head size.
\item Multiply per-token state by live tokens across requests.
\item Hybrid attention and recurrent state need architecture-specific accounting.
\item Use measured server capacity for the selected model and build.
\end{itemize}
\end{frame}
\begin{frame}[t,fragile]{Four sequences are not four full contexts}
\fontsize{11.5}{14}\selectfont
\begin{itemize}\setlength{\itemsep}{2mm}
\item The baseline allows up to four active sequences.
\item Its configured per-request limit is 262,144 tokens.
\item Four full requests would need state for 1,048,576 live tokens.
\item Scheduler admission and memory capacity must both be checked.
\end{itemize}
\vfill{\fontsize{6.5}{8}\selectfont\color{UTgray} Sources: user. See references and instructor notes.\par}
\end{frame}
\begin{frame}[t,fragile]{GPU topology affects distributed serving}
\fontsize{11.5}{14}\selectfont
\begin{itemize}\setlength{\itemsep}{2mm}
\item Tensor parallelism partitions operations within model layers.
\item Pipeline parallelism partitions layer groups.
\item Expert parallelism distributes MoE experts across devices.
\item Communication and per-device balance can limit scaling.
\end{itemize}
\vfill{\fontsize{6.5}{8}\selectfont\color{UTgray} Sources: parallel, glmrecipe. See references and instructor notes.\par}
\end{frame}
\begin{frame}[t,fragile]{Host offload exchanges capacity for other costs}
\fontsize{11.5}{14}\selectfont
\begin{itemize}\setlength{\itemsep}{2mm}
\item Host RAM can store weights that do not fit in VRAM.
\item CPU memory bandwidth, NUMA placement, and transfers become important.
\item An engine may transfer weights or compute selected work on CPU.
\item Measure usable interaction latency, not just successful loading.
\end{itemize}
\end{frame}
\begin{frame}[t,fragile]{Plan the complete serving host}
\fontsize{11.5}{14}\selectfont
\begin{itemize}\setlength{\itemsep}{2mm}
\item RAM and storage for model staging, runtime, logs, and rollback.
\item PCIe lanes and device topology appropriate to the GPU count.
\item Power, cooling, rack space, and maintenance access.
\item Network capacity and account controls for the class workload.
\end{itemize}
\end{frame}
\begin{frame}[t,fragile]{Read the primary settings before changing them}
\fontsize{10.5}{12.5}\selectfont\setlength{\tabcolsep}{4pt}\renewcommand{\arraystretch}{1.15}
\begin{tabular}{>{\raggedright\arraybackslash}p{\dimexpr 0.27\linewidth-4.32pt\relax}>{\raggedright\arraybackslash}p{\dimexpr 0.73\linewidth-11.68pt\relax}}
\toprule
\textbf{Setting} & \textbf{Supplied value} \\ \midrule
Weight checkpoint & unsloth/Qwen3.8-27B-NVFP4 \\[2pt]
Context / max sequences & 262144 / 4 \\[2pt]
GPU memory utilization & 0.92 \\[2pt]
KV cache mode & auto \\[2pt]
Speculation & MTP, 3 draft tokens \\[2pt]
Chat template & Local chat\_template.jinja \\[2pt]
\bottomrule\end{tabular}\par\vspace{2mm}
\vfill{\fontsize{6.5}{8}\selectfont\color{UTgray} Sources: user. See references and instructor notes.\par}
\end{frame}
\begin{frame}[t,fragile]{Tool use depends on the serving contract}
\fontsize{11.5}{14}\selectfont
\begin{itemize}\setlength{\itemsep}{2mm}
\item The baseline uses reasoning parser qwen3.
\item It enables automatic tool choice with parser qwen3\_coder.
\item The custom template enables and preserves thinking.
\item Test structured calls through OpenCode after any template change.
\end{itemize}
\vfill{\fontsize{6.5}{8}\selectfont\color{UTgray} Sources: user. See references and instructor notes.\par}
\end{frame}
\begin{frame}[t,fragile]{Keep client and server configuration distinct}
\fontsize{11.5}{14}\selectfont
\begin{itemize}\setlength{\itemsep}{2mm}
\item Client .env: API URL, issued key, served model, and limits.
\item Server: weight path, hardware, runtime, template, TLS, and admission.
\item The model name in the client must match the served alias.
\item Students do not need GPU access to use the private class endpoint.
\end{itemize}
\end{frame}
\begin{frame}[t,fragile]{Validate service readiness in stages}
\fontsize{11.5}{14}\selectfont
\begin{itemize}\setlength{\itemsep}{2mm}
\item Process: inspect startup errors and selected model settings.
\item Transport: verify TLS, authentication, and /v1/models.
\item Inference: complete a small public prompt.
\item Workflow: execute a local tool and verify its result through OpenCode.
\end{itemize}
\end{frame}
\begin{frame}[t,fragile]{Benchmark metrics answer different questions}
\fontsize{10.5}{12.5}\selectfont\setlength{\tabcolsep}{4pt}\renewcommand{\arraystretch}{1.15}
\begin{tabular}{>{\raggedright\arraybackslash}p{\dimexpr 0.27\linewidth-4.32pt\relax}>{\raggedright\arraybackslash}p{\dimexpr 0.73\linewidth-11.68pt\relax}}
\toprule
\textbf{Metric} & \textbf{Definition for this lab} \\ \midrule
Time to first token & Request start to first generated token \\[2pt]
Decode rate & Generated tokens per decode time interval \\[2pt]
End-to-end latency & Request start to final response \\[2pt]
Aggregate throughput & All completed output tokens / wall time \\[2pt]
Task success & Cases meeting the predeclared correctness test \\[2pt]
\bottomrule\end{tabular}\par\vspace{2mm}
\end{frame}
\begin{frame}[t,fragile]{Use a controlled benchmark matrix}
\fontsize{11.5}{14}\selectfont
\begin{itemize}\setlength{\itemsep}{2mm}
\item Hold tasks, output budget, sampling, and versions fixed.
\item Compare short and long prompts at one and several clients.
\item Record latency percentiles, throughput, errors, and GPU memory.
\item Repeat runs and retain correctness checks alongside speed.
\end{itemize}
\end{frame}
\begin{frame}[t,fragile]{Choose a system for a class workload}
\fontsize{11.5}{14}\selectfont
\begin{itemize}\setlength{\itemsep}{2mm}
\item Eight students need responsive tool-use sessions.
\item Most requests are short; two may carry large code excerpts.
\item Would you maximize model size, context, or total throughput?
\item Propose a queue/admission rule and a benchmark to justify it.
\end{itemize}
\end{frame}
\begin{frame}[t,fragile]{One worked approach: sizing for the class workload}
\fontsize{11.5}{14}\selectfont
\begin{itemize}\setlength{\itemsep}{2mm}
\item Targets: responsive short interactive requests for eight students, with two large code excerpts handled explicitly; model size is not the objective.
\item Admission rule: a queue with a per-request token budget; large excerpts receive reserved capacity or are chunked so long jobs do not block short ones.
\item Justifying benchmark: a short-prompt matrix at concurrency 1 and 8, plus a two-long-prompt contention run, recording TTFT, latency percentiles, errors, and peak memory.
\item Decision rule: accept the configuration if short-request TTFT stays within target while the long requests are served, or if the long requests need a separate reserved capacity.
\end{itemize}
\end{frame}
\begin{frame}[t,fragile]{GLM is a separate deployment experiment}
\fontsize{11.5}{14}\selectfont
\begin{itemize}\setlength{\itemsep}{2mm}
\item Use its own checkpoint, model-specific template, and parsers.
\item The current recipe gives FP8 deployment and hardware guidance.
\item Begin with modest context and concurrency, then profile upward.
\item Compare task quality and latency against the Qwen baseline.
\end{itemize}
\vfill{\fontsize{6.5}{8}\selectfont\color{UTgray} Sources: glmrecipe, glm. See references and instructor notes.\par}
\end{frame}
\begin{frame}[t,fragile]{Quantization is a quality experiment too}
\fontsize{11.5}{14}\selectfont
\begin{itemize}\setlength{\itemsep}{2mm}
\item Weight precision and cache precision are independent choices.
\item Check kernel support for the exact accelerator and model.
\item Compare numerical-task, citation, and tool-use outcomes.
\item Accept a faster configuration only within its demonstrated scope.
\end{itemize}
\end{frame}
\begin{frame}[t,fragile]{DFlash2 belongs in an isolated experiment}
\fontsize{11.5}{14}\selectfont
\begin{itemize}\setlength{\itemsep}{2mm}
\item Status: instructor reports current instability.
\item Retain the working NVFP4/MTP environment as the baseline.
\item Compare crashes, output validity, accepted draft work, and latency.
\item Promote only after repeatable correctness and stability evidence.
\end{itemize}
\vfill{\fontsize{6.5}{8}\selectfont\color{UTgray} Sources: user. See references and instructor notes.\par}
\end{frame}
\begin{frame}[t,fragile]{Operate with a rollback path}
\fontsize{11.5}{14}\selectfont
\begin{itemize}\setlength{\itemsep}{2mm}
\item Record driver, CUDA, engine, model revision, and template hash.
\item Retain a working environment and a known-good configuration.
\item Test updates on public tasks before returning the class service.
\item Monitor failures and capacity without exposing credentials.
\end{itemize}
\end{frame}
\begin{frame}[t,fragile]{Diagnose three different out-of-memory cases}
\fontsize{11.5}{14}\selectfont
\begin{itemize}\setlength{\itemsep}{2mm}
\item Failure while loading weights.
\item Failure during a very long prompt prefill.
\item Failure only when several sessions are active.
\item Which measurements and configuration changes fit each case?
\end{itemize}
\end{frame}
\begin{frame}[t,fragile]{One worked approach: diagnosing the three OOM cases}
\fontsize{11.5}{14}\selectfont
\begin{itemize}\setlength{\itemsep}{2mm}
\item OOM at weight loading: measure checkpoint size, precision, and per-device placement; remedies are a different quantization, expert placement, or more devices.
\item OOM during a long prefill: measure activation and workspace headroom against the batch token budget; remedies are a lower prompt limit or a lower batch-token budget for that request.
\item OOM only under concurrency: measure the live cache state and the admission limit; remedies are a per-request token budget and an explicit queue.
\item Each case points to a different measurement; a single universal "increase memory utilization" remedy does not fit all three.
\end{itemize}
\end{frame}
\begin{frame}[t,fragile]{Record a benchmark row completely}
\fontsize{10.5}{12.5}\selectfont\setlength{\tabcolsep}{4pt}\renewcommand{\arraystretch}{1.15}
\begin{tabular}{>{\raggedright\arraybackslash}p{\dimexpr 0.23\linewidth-3.68pt\relax}>{\raggedright\arraybackslash}p{\dimexpr 0.77\linewidth-12.32pt\relax}}
\toprule
\textbf{Field group} & \textbf{Required values} \\ \midrule
Identity & Model/checkpoint, engine/build, GPU topology, precision, template hash. \\[2pt]
Workload & Prompt tokens, generated tokens, concurrency, task ID, cold/warm cache. \\[2pt]
Timing & TTFT, request latency, measurement interval and client-side timeout. \\[2pt]
Outcomes & Success/failure, validity checks, peak observed memory and raw log path. \\[2pt]
\bottomrule\end{tabular}\par\vspace{2mm}
\end{frame}
\begin{frame}[t,fragile]{Calculate throughput without conflating latency}
\fontsize{11.5}{14}\selectfont
\begin{itemize}\setlength{\itemsep}{2mm}
\item Illustrative run: four requests each generate 600 tokens; the batch completes in 40 seconds.
\item Aggregate output throughput is 2,400/40 = 60 tokens/s. This does not imply each user received 60 tokens/s.
\item If each request had 5 s to first token and 35 s of decode, its approximate decode rate is 600/35 = 17.1 tokens/s.
\item Report completion order, individual timings and failures; do not infer latency percentiles from one batch.
\end{itemize}
\end{frame}
\begin{frame}[t,fragile]{Lab: weights to a measured endpoint}
\fontsize{11.5}{14}\selectfont
\begin{itemize}\setlength{\itemsep}{2mm}
\item Inspect and launch the sanitized primary configuration on an approved host.
\item Connect local OpenCode through the shared .env workflow.
\item Benchmark a controlled matrix and verify engineering task outcomes.
\item Compare GLM or a second precision; explore DFlash2 only optionally.
\end{itemize}
\end{frame}
\begin{frame}[t,fragile]{Sample submission: the benchmark lab}
\fontsize{11.5}{14}\selectfont
\begin{itemize}\setlength{\itemsep}{2mm}
\item Setup record: the sanitized launch configuration, exact checkpoint and engine build, GPU topology, template hash, and the .env client workflow used for connection.
\item Benchmark table: one complete row per matrix cell (model, precision, context, concurrency) with timings, outcomes, peak memory, and the raw-log path.
\item Correctness alongside speed: engineering task outcomes (for example, the heat-balance value against the analytical result) recorded per configuration.
\item Recommendation and limits: a justified workload recommendation, untested configurations stated as untested, and the GLM or second-precision comparison treated as a separate experiment.
\end{itemize}
\end{frame}
\begin{frame}[t,fragile]{Exit ticket: defend a capacity claim}
\fontsize{11.5}{14}\selectfont
\begin{itemize}\setlength{\itemsep}{2mm}
\item Which bytes are estimates, and which are measurements?
\item Which context and concurrency were actually tested?
\item What hardware or runtime assumption could break the claim?
\item What result would make you reject the configuration?
\end{itemize}
\end{frame}
\begin{frame}[t,fragile]{A complete answer: defending the capacity claim}
\fontsize{11.5}{14}\selectfont
\begin{itemize}\setlength{\itemsep}{2mm}
\item Estimates versus measurements: the weight bytes are a lower-bound estimate from the checkpoint size; the per-device memory and latency percentiles are measured on the tested configuration.
\item Tested scope: name the exact model revision, context lengths, and concurrency levels actually run; untested combinations are outside the claim.
\item Fragile assumption: the usable bytes per GPU after driver and runtime overhead, or the kernel support for the precision on this accelerator.
\item Rejection criteria: an OOM at the advertised concurrency, or a correctness failure on a held-out task, invalidates the configuration.
\end{itemize}
\end{frame}
\begin{frame}[t]{References 1}
\fontsize{9}{12}\selectfont\begin{itemize}\setlength{\itemsep}{3mm}
\item \textbf{curriculum}: User-supplied PHANES curriculum: mod/01 through mod/08/README.md.
\item \textbf{colors}: \href{https://umac.utexas.edu/brand-center/colors/}{umac.utexas.edu/brand-center/colors}
\item \textbf{user}: User-supplied vllm(1).zip: vllm/vllm-qwen38-nvfp4.sh, README.md, chat\_template.jinja; reviewed 2026-09-05.
\item \textbf{glm}: \href{https://huggingface.co/zai-org/GLM-5.3-Flash}{huggingface.co/zai-org/GLM-5.3-Flash}
\item \textbf{qmax}: \href{https://www.alibabacloud.com/en/press-room/alibaba-unveils-qwen3-8-max?_p_lc=1}{alibabacloud.com/en/press-room/alibaba-unveils-qwen3-8-max?\_p\_lc=1}
\item \textbf{rtx}: \href{https://www.nvidia.com/en-us/products/workstations/professional-desktop-gpus/rtx-pro-6000.md}{nvidia.com/en-us/products/workstations/professional-desktop-gpus/rtx-pro-6000.md}
\item \textbf{b300}: \href{https://www.nvidia.com/en-us/data-center/dgx-b300/}{nvidia.com/en-us/data-center/dgx-b300}
\end{itemize}\end{frame}
\begin{frame}[t]{References 2}
\fontsize{9}{12}\selectfont\begin{itemize}\setlength{\itemsep}{3mm}
\item \textbf{glmrecipe}: \href{https://recipes.vllm.ai/zai-org/GLM-5.3-Flash}{recipes.vllm.ai/zai-org/GLM-5.3-Flash}
\item \textbf{parallel}: \href{https://docs.vllm.ai/en/v0.5.1/serving/distributed_serving.html}{docs.vllm.ai/en/v0.5.1/serving/distributed\_serving.html}
\end{itemize}\end{frame}
\end{document}